\section{Introduction}
\label{sec:intro}

Risk-aware navigation has a selectivity gap. When a lower-risk feasible
maneuver exists, the robot should expose the evasive degree of freedom and
take it; when the apparent escape is blocked, corrupted, or not yet open, it
should suppress that same degree of freedom rather than hallucinate a detour.
Classical planners~\citep{hart1968formal,karaman2011sampling},
MPC/MPPI-style samplers~\citep{williams2017mppi}, reactive controllers, and
learned policies~\citep{schulman2017proximalpolicyoptimizationalgorithms,haarnoja2018softactorcriticoffpolicymaximum}
all address parts of this problem, but none make this
activation/suppression asymmetry an explicit structural property: risk costs
tend either to pull everywhere a gradient exists, causing false maneuvers, or
to be damped enough that the system misses real escape opportunities.

The issue is acute in off-road terrain and local traffic. A patch that looks
clear from geometry alone may be wet clay; a corridor that looks low-risk in
the semantic map may be physically blocked; a lane change may be useful behind
a slow leader but unsafe when the adjacent lane is boxed in. The decision
problem is therefore not just to reweight a scalar cost. The agent must decide
which force channels should be active under the current local context, and
which should remain latent.

Our key observation is that a single additive Hamiltonian energy term can make
this selectivity structural. Adding $\Ham_{\rm ctx}$ to a geometry-only
port-Hamiltonian policy simultaneously creates (i) a context force channel
in the momentum update, (ii) a parameter subspace and sensitivity stream for
learning when that channel matters, and (iii) a route by which CVaR gradients
from rare failures act directly on the closed-loop field.
Proposition~\ref{prop:enrichment} formalizes this coupling: enrichment is not merely a
new cost term, but a coordinated enlargement of the force field, learning
pathway, and update timescale.

The resulting policy remains a zero-replan local field. New sensory
information changes the cotangent update, not an external planner or critic.
A local affordance gate tests only the currently sensed patch: when a feasible
lower-risk primitive exists, the soft-risk force bends the field toward it;
when the candidate is blocked or insufficiently better, the soft channel is
suppressed and the rollout stays near the geometry-only policy.

\begin{figure*}[t]
    \centering
    \includegraphics[width=0.98\textwidth]{NeurIPS_2026/figures/teaser_decision_field.png}
    \caption{
    \textbf{Selective reshaping of the decision field.}
    \textbf{(A)} Geometrically feasible maneuvers can differ in material
    risk. \textbf{(B)} Adding
    \( -\tau \nabla_{\Pos} \Ham_{\mathrm{ctx}} \) to the cotangent update
    creates a context-force channel. \textbf{(C)} The channel bends toward
    a safer lane when one is feasible, but remains negligible when escape is
    boxed in. Sec.~\ref{sec:experiments} measures this activation/suppression
    signature directly.
    }
    \label{fig:teaser}
\end{figure*}

Prior Hamiltonian and port-Hamiltonian learners~\citep{greydanus2019hamiltonian,desai2021porthamiltonian,ellendula2025grlsnam}
provide structured vector fields, but their geometry-only policy is blind to
material/context risk. We extend that policy with soft-risk and hard-hazard
force channels and train with CVaR to emphasize rare consequential
transitions~\citep{chow2015risk}. The primary head-to-head benchmark is a
delayed local escape: the lateral maneuver is blocked early, necessary later,
and must be timed without global replanning.

% A prior geometry-only Hamiltonian formulation~\citep{ellendula2025grlsnam} encodes goal
% attraction, obstacle repulsion, and
% dissipation, and outperforms reactive planners and deep RL baselines
% under matched local sensing. But geometry-only structure is materially
% blind: two trajectories with identical clearance receive identical
% forces, even when one crosses water and the other crosses pavement.
% This paper extends the formulation to the material-aware
% risk-sensitive setting through three coordinated additions: (i) a
% soft-risk potential $\Ham_{\mathrm{ctx}}^{\mathrm{soft}}$ and hard
% hazard barrier $\Ham_{\mathrm{haz}}^{\mathrm{hard}}$ that introduce
% material force channels directly into the $\Mom$-update; (ii) a CVaR
% tail-risk objective that drives episodic gradient updates against
% catastrophic outcomes rather than expected
% cost~\citep{chow2015risk}; and (iii) a three-loop adaptation
% hierarchy operating at segment, episode, and curriculum timescales.
% These additions are not independent. We show that they are linked by
% a \emph{learnable energy enrichment principle}: a single new additive
% energy term in $\Ham$ induces a corresponding force channel, parameter
% subspace, gradient pathway, and update timescale. The context-enriched learner is
% therefore a
% structurally consistent enlargement of the geometry-only learner
% along four coupled axes. The resulting field can be evaluated through
% selectivity metrics: material force channels should be active when the local
% geometry admits a lower-risk alternative and negligible otherwise.
% In this sense, the context-enriched field learns not only how much to value each risk term,
% but which context-conditioned degrees of freedom the closed-loop system
% should expose.

% We validate this with three kinds of evidence. \emph{Static
% semantic-risk reduction:} on the DFC2018-derived material-aware
% benchmark, context-enriched field substantially improves the learned geometry-only
% scaffold, reducing hard-hazard length, cumulative risk exposure, and
% oscillation while preserving path efficiency. We also compare against
% representative risk-aware planners and an oracle reference, and report
% where these planners remain stronger on raw static-map risk. \emph{Dynamic
% context-sensitive behavior:} in highway-env scenarios with moving
% traffic, context-enriched field produces lane changes when a slow leader admits a
% safer adjacent lane, removes default off-road failures, and remains
% conservative when no adjacent lane is available. \emph{Mechanistic
% agreement:} the per-step force decomposition shows that the lateral
% component of the material force is selective in exactly the way the
% behavior is -- large in the escape-available regime, near-zero in the
% escape-blocked regime, and near-zero on benign default scenes.

Our contributions are:
\begin{itemize}
  \item \textbf{Selective risk-field enrichment.} One additive
        context-energy term creates a force channel that activates
        when the local scene admits a lower-risk feasible maneuver and
        suppresses when it does not. This turns risk adaptation into a
        context-conditioned change in the closed-loop field, rather
        than a uniform reweighting of a scalar cost.

  \item \textbf{A learnable energy enrichment principle.} We prove
        (Proposition~\ref{prop:enrichment}) that, within the class of
        port-Hamiltonian learners considered here, a new additive
        energy term induces a new force channel, parameter subspace,
        sensitivity stream, and update timescale. The context-enriched
        enlargement is structurally consistent rather than ad hoc.

  \item \textbf{Falsifiable evaluation across four settings.} We lead with a
        delayed-required-escape benchmark where the method reduces premature
        activation while increasing success with zero replans, then measure
        spatial selectivity (CAR/FAR/SR), temporal selectivity (false
        pre-activation, suppress rate, reaction delay), and empirical
        tail-risk reduction (violation CVaR) across RELLIS-3D, RELLIS-Dyn,
        DFC2018, and highway-env.
\end{itemize}

The central message of the paper is not that the learner becomes
risk-aware. It is that the learner becomes risk-aware
\emph{selectively}: reshaping the decision field exactly when the
geometry of the situation permits a safer alternative, exposing only the
useful degrees of freedom, and leaving the geometry-only behavior unchanged otherwise.




% \section{Introduction}
% \label{sec:intro}

% Consider an agent navigating terrain whose risk structure is revealed
% only locally: a stretch of road that looks identical from above may be
% asphalt or wet clay, a corridor between two obstacles may be a safe
% shortcut or a soft shoulder that fails under load. Two maneuvers that
% are geometrically feasible — both clear of obstacles, both reaching
% the goal — are not equally good when one carries higher latent
% material risk. The agent's task is not to find a feasible path once,
% but to \emph{continually shift its behavior toward lower-risk feasible
% paths} as material structure is revealed through experience, and to
% do so under the same local-information regime in which it must
% already act.

% Existing methods address parts of this problem but each leaves
% something fixed. Classical planners~\citep{hart1968formal,karaman2011sampling}
% and risk-aware MPC~\citep{williams2017mppi} adapt the \emph{path}
% against a given cost map. Deep reinforcement
% learning~\citep{schulman2017proximalpolicyoptimizationalgorithms,
% haarnoja2018softactorcriticoffpolicymaximum} adapts the \emph{policy}
% weights against a fixed reward. Safe and risk-sensitive
% RL~\citep{achiam2017constrained,chow2015risk} adapt the
% \emph{constraint set} or the \emph{risk-weighted cost}. In each case,
% the closed-loop decision field — the vector field that actually
% determines how the agent moves at every state — is typically reshaped
% only indirectly, through an external planner, critic, or cost
% evaluator that sits outside the dynamics. The decision field itself
% is not the object being learned.

% We take the opposite view. We represent the policy as the time-$\Delta
% t$ flow map of a reduced port-Hamiltonian energy
% $\Ham : T^*\mathcal{Q} \to \mathbb{R}$, so that the decision field
% \emph{is} the Hamiltonian vector field, and adapting behavior
% \emph{is} reshaping that field. The cotangent variable $\Mom$ — the
% generalized momentum — acts as an internal evaluator, predictor, and
% guide: new sensory information enters by modifying $\Mom$, which in
% turn shapes the force balance that drives the next configuration
% update. There is no external critic and no external replanner; the
% risk assessment enters through the same dynamics that generate motion.

% \begin{figure*}[t]
%     \centering
%     \includegraphics[width=0.98\textwidth]{Master_s_Thesis/figs/teaser_decision_field.png}
%     \caption{
%     \textbf{Reshaping the decision field under local material risk.}
%     \textbf{(A)} Two maneuvers can be geometrically feasible yet materially different in risk.
%     \textbf{(B)} Our method injects local material risk directly into the cotangent update through the added term
%     \( -\tau \nabla_{\Pos} \Ham_{\mathrm{ctx}} \), thereby reshaping the decision field itself.
%     \textbf{(C)} Under the same local information, this yields a lower-risk feasible path than a material-blind baseline, with comparable path efficiency and fewer oscillations.
%     }
%     \label{fig:teaser}
% \end{figure*}

% In the geometry-only GRL-SNAM
% formulation~\citep{ellendula2025grlsnam}, this field encodes goal
% attraction, obstacle repulsion, and dissipation, and the resulting
% learner outperforms reactive planners and deep RL baselines under
% matched local sensing. But geometry-only structure is materially
% blind: two trajectories with identical clearance receive identical
% forces, even when one crosses water and the other crosses pavement.
% The learner has no force channel through which material risk can act
% on the dynamics, no objective that concentrates gradient on rare
% catastrophic outcomes, and no parameter subspace in which
% material-specific corrections can be stored.

% This paper extends the formulation to the material-aware
% risk-sensitive setting through three coordinated additions: (i) a
% soft-risk potential $\Ham_{\mathrm{ctx}}^{\mathrm{soft}}$ and hard
% hazard barrier $\Ham_{\mathrm{haz}}^{\mathrm{hard}}$ that introduce
% material force channels directly into the $\Mom$-update; (ii) a CVaR
% tail-risk objective that drives episodic gradient updates against
% catastrophic outcomes rather than expected
% cost~\citep{chow2015risk}; and (iii) a three-loop adaptation
% hierarchy operating at segment, episode, and curriculum timescales.
% These additions are not independent. We show that they are linked by
% a \emph{learnable energy enrichment principle}: a single new additive
% energy term in $\Ham$ induces a corresponding force channel,
% parameter subspace, gradient pathway, and update timescale. The
% resulting context-enriched learner is therefore not
% the geometry-only learner with risk costs bolted on — it is a
% structurally consistent enlargement of the geometry-only learner
% along four coupled axes.
% Figure~\ref{fig:intro:internal_evaluator} makes the contrast
% concrete: the $\Mom$-update gains exactly one additive term, and
% every other change in the algorithm follows from it.

% Our contributions are:
% \begin{itemize}
%   \item \textbf{Material-aware continual Hamiltonian learning under
%         local information.} A structured navigation method in which
%         the policy is a port-Hamiltonian phase-space flow continually
%         reshaped by material-aware force channels, with $\Mom$ as the
%         internal evaluator rather than an external critic or
%         replanner.
%   \item \textbf{Risk-improving field adaptation, not just
%         risk-weighted cost.} The added force channels and CVaR-driven
%         episodic updates change the induced vector field so that the
%         learner continually shifts toward lower-risk feasible
%         trajectories under the same local geometric scaffold — a
%         behavioral, not just representational, change.
%   \item \textbf{A learnable energy enrichment principle.} We prove
%         (Theorem~\ref{thm:enrichment}) that a new additive energy
%         term induces a new force channel, parameter subspace,
%         sensitivity stream, and update timescale, making the context-enriched field
%         enlargement structurally consistent rather than ad hoc.
%   \item \textbf{Empirical gains under matched local information.} On
%         DFC2018-derived material-aware navigation, the proposed
%         method reduces cumulative risk by $56.4\%$ over the
%         geometry-only baseline, outperforms same-information local
%         baselines (risk-weighted DWA, CBF+risk cost, local risk-A*),
%         generalizes to the full held-out test set, and degrades
%         gracefully under noisy risk maps.
% \end{itemize}

% The central message of the paper is not that the learner becomes
% material-aware. It is that material-awareness emerges from a
% continually reshaped internal decision field, and that this
% reshaping is the right object of learning when the environment's
% risk structure is revealed only through experience.
